\question{5.m5}{Van een normaal verdeelde kansvariabele $X$ is gegeven dat deze een standaarddeviatie heeft die gelijk is aan $9$.
Verder is gegeven dat $P(X < 15.5) = 0.3085$.
De waarde van $\mu$ van deze variabele is dan $\ldots$}
\begin{enumerate}[label=(\alph*)]
    \item $16$.
    \item $17$.
    \item $20$.
    \item $11$.
\end{enumerate}
\answer{
    In dit geval willen we bepalen wat de waarde van $\mu$ is zodanig dat geldt:
    \[
        \normalcdf(\lb=-\GRinfinity, \ub=15.5, \mu=?, \sigma=9) = 0.3085.
    \]
    In de grafische rekenmachine kun je de functie \enquote{numeric solver} gebruiken (TI84-Plus CE: MATH-C en Casio fx-CG50: OPTN-CALC-SolveN) en het volgende invullen:
    \begin{align*}
        E_1 &= \normalcdf(-\GRinfinity, 15.5, X, 9) \\
        E_2 &= 0.3085
    \end{align*}
    Zodra je dit hebt ingevoerd en op OK hebt gedrukt (via de knop \enquote{graph}), krijg je een scherm te zien met $X=\ldots$ en $\text{bounds}=\ldots$.
    
    De $X=$ geeft een eerste gok aan voor de waarde van $X$. Voor nu maakt het niet uit wat je hier invult, dus stel $X=0$.
    De bounds geven aan in welk gebied de rekenmachine moet gaan zoeken, hier kun je het beste even de standaardinstellingen gebruiken.
    Als je met de cursor op de regel staat met $X=\ldots$, dan komt er SOLVE in beeld te staan.
    Druk nu op SOLVE (nogmaals de knop \enquote{graph}).
    De rekenmachine gaat nu de vergelijking $E_1 = E_2$ oplossen, waaruit een oplossing $X$ komt.

    Dit geeft in ons geval een snijpunt voor $X = 20.0010$.
    Dit betekent dat $\mu$ een waarde van ongeveer $20$ heeft.
    Het juiste antwoord is dus (c).

    {
        \noindent \textit{\textbf{Noot:} bij Casio fx-CG50 is de syntax gelijk aan $\text{SolveN}(\text{NormCD}(-\GRinfinity, 15.5, 9, x)=0.3085, x)$.}
    }
}